%
% Capítulo 3
%
\chapter{Arquitetura da Solução} \label{cap:arquitetura}

A solução desenvolvida assenta numa arquitetura composta por dois blocos principais: uma aplicação móvel Android,
responsável pela interação com os utilizadores, e um backend com API REST, responsável pela persistência central
de dados e pela execução da lógica associada à gestão de informação do sistema. Esta separação permite distribuir
responsabilidades de forma clara entre os componentes, simplificando a evolução da solução e favorecendo a manutenção do código.

De forma geral, a aplicação Android permite ao utilizador interagir com o sistema de acordo com o seu perfil,
nomeadamente professor ou aluno, enquanto o backend disponibiliza os serviços necessários para registo e consulta
de informação. A comunicação entre ambos é realizada através de pedidos HTTP a uma API REST, permitindo que o
frontend envie dados estruturados e obtenha respostas consistentes para os vários fluxos funcionais da aplicação.

\section{Visão Global da Arquitetura}

A arquitetura adotada segue uma abordagem em camadas, onde cada componente desempenha uma função específica.
No frontend, a interface foi construída com Jetpack Compose, recorrendo a \textit{ViewModels} e repositórios
para organizar o estado da aplicação e a lógica de acesso a dados.
No backend, a API foi implementada em Ktor, com persistência em PostgreSQL e uma camada de acesso à base de dados
suportada por Exposed e HikariCP.

Em termos funcionais, a arquitetura pode ser vista como composta pelos seguintes elementos:
\begin{itemize}
    \item aplicação Android, que apresenta os ecrãs e gere a interação com o utilizador;
    \item camada de lógica no frontend, responsável por coordenar o estado da interface e os fluxos de utilização;
    \item API REST no backend, que centraliza as operações de leitura, escrita e consulta de dados;
    \item base de dados PostgreSQL, que armazena a informação persistente do sistema;
    \item armazenamento local com Room, utilizado como suporte local no dispositivo Android.
\end{itemize}

Esta organização permite que a aplicação continue a evoluir de forma modular, com fronteiras relativamente claras entre
apresentação, lógica de negócio, comunicação remota e persistência.

\section{Arquitetura do Frontend}

O frontend da solução foi desenvolvido em Kotlin para Android, utilizando Jetpack Compose como tecnologia principal para
construção da interface. Esta escolha permite uma abordagem declarativa ao desenho dos ecrãs, facilitando a gestão de
estados e a criação de componentes reutilizáveis.

A aplicação encontra-se organizada em várias camadas, com destaque para:
\begin{itemize}
    \item \textit{screens}, responsáveis pela apresentação visual das funcionalidades;
    \item \textit{ViewModels}, responsáveis pela gestão do estado e coordenação da lógica de apresentação;
    \item \textit{repositories}, responsáveis pelo acesso a dados locais e remotos;
    \item entidades locais em Room, que suportam persistência local no dispositivo.
\end{itemize}

Do ponto de vista funcional, a aplicação distingue dois perfis de utilizador: professor e aluno. Cada perfil possui ecrãs e ações específicas.
O professor pode gerir a sua disponibilidade, definir restrições e consultar os alunos associados. O aluno pode escolher um professor e registar
a sua disponibilidade. Esta separação por papéis foi integrada tanto na navegação da aplicação como na própria lógica dos \textit{ViewModels}.

\section{Arquitetura do Backend}

O backend foi desenvolvido em Kotlin com Ktor, funcionando como servidor de aplicação e ponto central de integração de dados.
A API implementada segue o estilo REST e expõe endpoints para operações como:

\begin {itemize}
    \item criação e consulta de professores e alunos;
    \item associação de alunos a professores;
    \item registo e consulta de disponibilidades;
    \item registo e consulta de restrições;
    \item criação de propostas de horário.
\end{itemize}

A estrutura interna do backend encontra-se organizada em componentes como:
\begin{itemize}
    \item definição de rotas, responsável por mapear pedidos HTTP para operações concretas;
    \item modelos e DTOs, que representam os dados trocados entre cliente e servidor;
    \item serviços, que encapsulam partes relevantes da lógica do sistema;
    \item acesso à base de dados, com apoio de Exposed para modelação e manipulação de tabelas.
\end{itemize}

A escolha de Ktor permitiu construir um backend leve, relativamente simples de estruturar e bem alinhado com o restante ecossistema Kotlin do projeto.
Esta solução revelou-se adequada para suportar o conjunto de operações necessárias à versão beta e para validar a integração com a aplicação móvel.

\section{Persistência de Dados}

A persistência do sistema está dividida em dois níveis complementares. No backend, os dados principais são armazenados em PostgreSQL,
funcionando esta base de dados como repositório central do sistema. No frontend, foi utilizado Room para persistência local,
permitindo manter dados no dispositivo e suportar alguns fluxos locais ou híbridos durante a evolução da aplicação.

Na base de dados central são armazenadas entidades essenciais ao domínio do problema, tais como:
\begin{itemize}
    \item professores;
    \item alunos;
    \item disponibilidades;
    \item restrições.
\end{itemize}

A utilização de Room no frontend permite preservar alguma informação localmente, o que facilita o funcionamento de certas operações
e reduz dependência imediata da comunicação remota em todos os momentos. No entanto, o objetivo global da arquitetura é que a API
e a base de dados remota assumam progressivamente o papel de fonte principal de verdade do sistema.

\section{Comunicação entre Frontend e Backend}

A comunicação entre a aplicação Android e o backend é feita através de pedidos HTTP para a API REST, utilizando o endereço local do
servidor durante o desenvolvimento. Os dados trocados entre os componentes são serializados em formato JSON, permitindo uma
representação simples e facilmente interpretável por ambas as partes.

Esta comunicação suporta já um conjunto importante de operações, incluindo:
\begin{itemize}
    \item criação de professores e alunos;
    \item consulta de listas de utilizadores;
    \item associação de alunos a professores;
    \item registo e consulta de disponibilidades;
    \item registo e consulta de restrições;
    \item criação de propostas de horário.
\end{itemize}

A existência desta comunicação remota permite validar o sistema como uma solução efetivamente distribuída, em que a aplicação cliente
e o backend colaboram para concretizar os fluxos principais do projeto. A versão beta representa uma etapa importante desta integração,
demonstrando que os componentes já comunicam entre si e que a base funcional do sistema se encontra operacional.

\section{Considerações sobre a Arquitetura Atual}

A arquitetura adotada foi escolhida tendo em conta a necessidade de construir uma solução progressivamente, validando primeiro os fluxos
principais e a integração entre componentes, antes de avançar para otimizações mais profundas. A existência simultânea de persistência local
e remota reflete a evolução incremental do projeto, permitindo apoiar o desenvolvimento da aplicação móvel ao mesmo tempo que se consolida a
API e a base de dados central.

Embora a arquitetura atual já permita demonstrar de forma consistente o funcionamento da solução, existem ainda aspetos a consolidar,
nomeadamente o reforço da sincronização entre dados locais e remotos e a clarificação do papel de cada camada em alguns fluxos específicos.
Ainda assim, a base estrutural da solução encontra-se definida e suficientemente sólida para suportar a continuação do desenvolvimento até à versão final.
